\section{Project Controls} \label{sec:control}
DM operations tracks work and milestones in Jira following \citeds{RTN-005}.

Risk management follows the Operations risk management plan \citedsp{rdo-71}.

Security is covered in \cite{RTN-030}.

Disaster recovery is covered in various documents:

\begin{itemize}
\item \citeds{ITTN-058} covers Disaster Recovery for Infrastructure Support Devices
\item \citeds{ITTN-057} covers Disaster Recovery for Computing
\item \citeds{ITTN-056} covers Disaster Recovery for the  Network
\end{itemize}


\subsection{DM Change Control Board \label{sect:dmccb}}

The DMCCB has responsibility for issues similar to those of the Rubin Change Control Board, but focused on the DM Subsystem.
The DMCCB reviews and approves changes to all baselines in the Subsystem, including proposed changes to any change controlled documents.
This is done via the RFC mechanism.
The Technical Baseline, including software/hardware and documentation, is produced by DM and controlled by the DMCCB.
DMCCB validates that the form and content of the Technical Baseline is consistent with Rubin project standards.
See also \url{https://rubinobs.atlassian.net/wiki/spaces/DM/pages/48828899/DMCCB}

\begin{itemize}
\item Responsibilities:
        \begin{itemize}
        \item Review and approve/reject proposed changes to baselined items - any Rubin Change Request must go through the DMCCB before being submitted to the project CCB.
        \item Review all RFCs and approves \textit{flagged} RFCs prior to 'Adoption'
        \item Monitor and approve DM software releases
        \item Monitor the status of issues in the DM project on Jira
        \item Ensure that the DM follows Rubin and DM configuration control processes.
        \end{itemize}
\item Membership:
        \begin{itemize}
        \item Core members:
                \begin{itemize}
                \item Chair Tim Jenness.
                \item DMAD
                \item DMSS
                \item USDF representative
                \item Release Manager
                \end{itemize}
        \item Optional members (required when topics to discuss are relevant to their areas of expertise):
                \begin{itemize}
                \item Deputy DMSS, when DMS is not available
                \item Team Leads, who can delegate as needed % e.g.Jim Bosch for pipelines
		\item Any team member who wishes to discuss an RFC or ticket.
                \end{itemize}
        \item For on-line virtual meetings, if a consensus or quorum is not reached within one week, the DMAD will make a unilateral decision
        \item DMAD can also make unilateral decisions in cases of urgency. In that case DMCCB will assess the change \textit{a posteriori}.
        \end{itemize}
\end{itemize}

The DMCCB will meet virtually, every week for 30 minutes. Agenda will be available beforehand.
Urgent decisions can be taken offline, outside the weekly meeting, in a modality to be defined by the DMCCB itself (email or slack channel).

All RFCs that imply one of the following changes:

\begin{itemize}
\item Changes to controlled documents
\item API changes to the codebase, including deprecation
\item Data model changes
\end{itemize}

need to be \textit{flagged} and therefore approved by the DMCCB, as detailed in the \href{https://developer.lsst.io/communications/rfc.html#rfc-exceptions}{Developer Guide}.




\subsection{Work Breakdown Structure} \label{sec:wbs}
Table \tabref{tab:dmwbs} gives the WBS structure for DM in operations.

\normalsize \begin{longtable} {|p{0.1\textwidth}|l|l|} \caption{WBS elements for Rubin Observatory Data Management Operations \label{tab:dmwbs}}\\ 
\hline 
\textbf{WBS}&\textbf{Department}&\textbf{Team} \\ \hline
{3.1}&{Data Management Operations}&{RDM Management} \\ \hline
{3.10}&{Data Management Operations}&{Service Quality and Reliability Engineering} \\ \hline
{3.10.1}&{Data Management Operations}&{Service Quality and Reliability Engineering} \\ \hline
{3.10.2}&{Data Management Operations}&{Service Quality and Reliability Engineering} \\ \hline
{3.10.3}&{Data Management Operations}&{Service Quality and Reliability Engineering} \\ \hline
{3.11}&{Data Management Operations}&{Qserv} \\ \hline
{3.11.1}&{Data Management Operations}&{Qserv} \\ \hline
{3.12}&{Data Management Operations}&{Build Engineering} \\ \hline
{3.12.1}&{Data Management Operations}&{Build Engineering} \\ \hline
{3.2}&{Data Management Operations}&{US Data Facility} \\ \hline
{3.2.1}&{Data Management Operations}&{US Data Facility} \\ \hline
{3.2.2}&{Data Management Operations}&{US Data Facility} \\ \hline
{3.2.3}&{Data Management Operations}&{US Data Facility} \\ \hline
{3.2.4}&{Data Management Operations}&{US Data Facility} \\ \hline
{3.2.6}&{Data Management Operations}&{US Data Facility} \\ \hline
{3.2.7}&{Data Management Operations}&{US Data Facility} \\ \hline
{3.3}&{Data Management Operations}&{France Data Facility} \\ \hline
{3.4}&{Data Management Operations}&{UK Data Facility} \\ \hline
{3.5}&{Data Management Operations}&{Base \& Summit Data Facilities} \\ \hline
{3.5.1}&{Data Management Operations}&{Base \& Summit Data Facilities} \\ \hline
{3.6}&{Data Management Operations}&{Pipeline Middleware} \\ \hline
{3.6.1}&{Data Management Operations}&{Pipeline Middleware} \\ \hline
{3.7}&{Data Management Operations}&{Data Engineering} \\ \hline
{3.7.1}&{Data Management Operations}&{Data Engineering} \\ \hline
{3.7.2}&{Data Management Operations}&{Data Engineering} \\ \hline
{3.7.3}&{Data Management Operations}&{Data Engineering} \\ \hline
{3.8}&{Data Management Operations}&{Algorithms \& Pipelines} \\ \hline
{3.8.1}&{Data Management Operations}&{Algorithms \& Pipelines} \\ \hline
{3.8.2}&{Data Management Operations}&{Algorithms \& Pipelines} \\ \hline
{3.8.3}&{Data Management Operations}&{Algorithms \& Pipelines} \\ \hline
{3.9}&{Data Management Operations}&{Campaign Management} \\ \hline
{3.9.1}&{Data Management Operations}&{Campaign Management} \\ \hline
\end{longtable} \normalsize


